\section*{Chapter \ref{ch:m2:central-banks-and-the-short-rate} --- Central Banks and the Short Rate}

\begin{solution}{exo:m2:central-banks-and-the-short-rate:1}
Thursday accrues 1 day, Friday 3, Monday 1: $N = 5$.
$(1 + 0.0390/360)(1 + 0.0388 \times 3/360)(1 + 0.0392/360) - 1$, times
$360/5$, gives 3.8926\%. The day-weighted average is $(3.90 + 3 \times 3.88 +
3.92)/5 = 3.8920\%$; compounding adds 0.06 basis points over five days.
\end{solution}

\begin{solution}{exo:m2:central-banks-and-the-short-rate:2}
Corridor $2.90 - 2.50 = 40$ basis points; the refinancing rate is 15 basis
points above the deposit rate. At 2.440\% the benchmark is 6 basis points
below the deposit rate, which a bank with access to the facility would never
accept as a lender. But the rate measures what banks \emph{pay} to borrow,
including from money-market funds, insurers and other institutions that
cannot deposit at the central bank; the bound binds only lenders with access.
\end{solution}

\begin{solution}{exo:m2:central-banks-and-the-short-rate:3}
Central bank: securities $-10$, \omterm{def:m2:central-banks-and-the-short-rate:reserves}{reserves} $-10$. Insurer's bank: \omterm{def:m2:central-banks-and-the-short-rate:reserves}{reserves}
$-10$, insurer's deposit $-10$. Insurer: bonds $+10$, deposit $-10$. \omterm{def:m2:central-banks-and-the-short-rate:reserves}{Reserves}
and deposits are destroyed together; this is the mechanics of quantitative
tightening by sales.
\end{solution}

\begin{solution}{exo:m2:central-banks-and-the-short-rate:4}
Approximation: $r^2(N-1)/(2B) = 0.04^2 \times 90/720 = 2.00$ basis points.
Exactly: $[(1 + 0.04/360)^{91} - 1] \times 360/91 - 0.04 = 2.01$ basis points.
Over a quarter the term is not negligible for a swap of USD~1 billion: 2
basis points for a quarter is about USD~50\,000.
\end{solution}

\begin{solution}{exo:m2:central-banks-and-the-short-rate:5}
\omterm{def:m2:central-banks-and-the-short-rate:reserves}{Reserves} fall by USD~150 billion
(\cref{prop:m2:central-banks-and-the-short-rate:total}). In a corridor,
supply moves left along the steep part of the demand curve and the overnight
rate rises towards the lending rate unless the desk adds USD~150 billion,
by repo or purchases, the same morning; it must forecast the tax flow to do
so. In a floor, if \omterm{def:m2:central-banks-and-the-short-rate:reserves}{reserves} remain on the flat part the rate does not move and
the central bank does nothing. The size of the buffer it keeps is the price of
not having to forecast.
\end{solution}

\begin{solution}{exo:m2:central-banks-and-the-short-rate:6}
All-in $3.80 + 0.26161 + 1.50 = 5.56161\%$. Interest $50\,000\,000 \times
0.0556161 \times 91/360 = \text{USD}~702\,926$. The 26.161 basis points stand
for the bank credit and term premium that LIBOR contained and SOFR does not.
\end{solution}

\begin{solution}{exo:m2:central-banks-and-the-short-rate:7}
Lockout of two days: 3.7421\%, 0.27 basis points below the plain 3.7448\%
(30 September's quarter-end 3.95\% is replaced by 28 September's 3.87\%, and
29 September's by the same). Observation shift of five days: 3.6836\%, 6.12
basis points below, exactly the lookback's number. Both use the fixings of 25
August to 23 September; they differ only in which weight each fixing carries,
and here the two weight lists are the same multiset for each rate: the
four-day weight of the Labor Day weekend falls on a 3.62\% fixing in both
(28 August under the lookback, 4 September under the shift), and the
three-day weights likewise. With only two rate levels in the window, the
product of \cref{def:m2:central-banks-and-the-short-rate:arrears} is
identical. A hike inside the long weekend would break the tie.
\end{solution}

\begin{solution}{exo:m2:central-banks-and-the-short-rate:8}
In 2026 the Fed ran a floor with ample \omterm{def:m2:central-banks-and-the-short-rate:reserves}{reserves}: it raised the rate by
raising the rates it administers (interest on \omterm{def:m2:central-banks-and-the-short-rate:reserves}{reserves}, the reverse repo and
standing repo rates, the discount rate) and did not need to change the
quantity of \omterm{def:m2:central-banks-and-the-short-rate:reserves}{reserves} at all; its bill purchases in the same weeks
\emph{added} \omterm{def:m2:central-banks-and-the-short-rate:reserves}{reserves}. The statement describes a \omterm{def:m2:central-banks-and-the-short-rate:corridor}{corridor system} with scarce
\omterm{def:m2:central-banks-and-the-short-rate:reserves}{reserves}, where the desk drains or adds \omterm{def:m2:central-banks-and-the-short-rate:reserves}{reserves} each day to move the rate
along the steep part of the demand curve, roughly how the Fed operated before
2008.
\end{solution}

\begin{solution}{pb:m2:central-banks-and-the-short-rate:1}
\textbf{1.} $2 \times 10^9 \times (0.0225 \times 15 + 0.0250 \times 15)/360 =
\text{EUR}~3\,958\,333$.
\textbf{2.} $2 \times 10^9 \times (0.0265 \times 15 + 0.0290 \times 15)/360 =
\text{EUR}~4\,625\,000$.
\textbf{3.} EUR~666\,667, to the central bank, which pays the deposit rate
and receives the lending rate on the same \omterm{def:m2:central-banks-and-the-short-rate:reserves}{reserves}.
\textbf{4.} So that banks prefer to trade with each other: the gap rewards an
interbank market, which spreads information about banks and does the
central bank's allocating for it, and it makes the lending facility a
backstop rather than a habit.
\textbf{5.} Every rate in the market lies inside the corridor
(\cref{prop:m2:central-banks-and-the-short-rate:bounds}); moving both walls
moves the whole interval, and banks' relative bargaining is unchanged.
\textbf{6.} Middle: 2.45\% then 2.70\%. Interest EUR~4\,291\,667.
\textbf{7.} A gains EUR~333\,333 over the deposit facility; B saves
EUR~333\,333 over the lending facility.
\textbf{8.} The scarcity of \omterm{def:m2:central-banks-and-the-short-rate:reserves}{reserves}, and so each side's alternative:
with more excess \omterm{def:m2:central-banks-and-the-short-rate:reserves}{reserves} in the system lenders compete and the rate falls
towards the deposit rate; with fewer, borrowers compete and it rises. This is
the demand curve of \cref{fig:m2:central-banks-and-the-short-rate:corridor}.
\textbf{9.} 5 basis points on EUR~2 billion for 30 days is EUR~83\,333: A
gains EUR~416\,667, B saves EUR~250\,000. The premium is the price of B's
credit, which the facility does not charge because it takes collateral.
\textbf{10.} An unsecured benchmark (\texteuro STR, SONIA) measures what A
charges B without collateral; a secured one (SOFR) measures the repo market,
where collateral removes most of the credit premium.
\textbf{11.} To the deposit rate: with EUR~50 billion added, B has \omterm{def:m2:central-banks-and-the-short-rate:reserves}{reserves}
of its own, nobody needs to borrow, and a lender's only alternative is the
deposit facility.
\textbf{12.} Nothing: B does not borrow, or would pay no more than A earns at
the facility.
\textbf{13.} It raises the deposit rate (and the lending rate with it); it no
longer needs to adjust the quantity of \omterm{def:m2:central-banks-and-the-short-rate:reserves}{reserves}.
\textbf{14.} Not for the rate. Its balance sheet is now a separate
instrument: for longer-term yields, or simply a buffer that keeps \omterm{def:m2:central-banks-and-the-short-rate:reserves}{reserves} on
the flat part of the curve whatever the autonomous factors do.
\textbf{15.} $10^9 \times 0.0006 \times 30/360 = \text{EUR}~50\,000$: A earns
the gap between the fund's rate and the facility, without risk, because it
can use the facility and the fund cannot.
\textbf{16.} A smaller balance sheet and less interest paid to banks on
\omterm{def:m2:central-banks-and-the-short-rate:reserves}{reserves}; an active interbank market in which banks monitor each other and
the rate carries information.
\textbf{17.} Control of the rate without forecasting autonomous factors every
day; the freedom to use the balance sheet (purchases after a crisis, or
liquidity) without losing the rate.
\textbf{18.} With refinancing only 15 basis points above the deposit rate,
banks short of \omterm{def:m2:central-banks-and-the-short-rate:reserves}{reserves} borrow from the central bank cheaply, so \omterm{def:m2:central-banks-and-the-short-rate:reserves}{reserves} are
supplied on demand and the market rate stays close to the deposit rate even
as the balance sheet shrinks: a floor, supplied by demand.
\textbf{19.} Named result: \emph{the gain from trade in \omterm{def:m2:central-banks-and-the-short-rate:reserves}{reserves}} over the
period is $(i_L - i_D) \times \text{amount} \times \text{days}/360 = 0.40\%
\times \text{EUR}~2\,\text{billion} \times 30/360 = \text{EUR}~666\,667$,
shared between A and B by where the interbank rate settles. In a floor it
is zero: no bank is short, and the rate sits on the deposit rate.
\textbf{20.} The rates the central bank administers, and where the supply of
\omterm{def:m2:central-banks-and-the-short-rate:reserves}{reserves} it chooses meets banks' demand for them.
\end{solution}

\begin{solution}{iq:m2:central-banks-and-the-short-rate:1}
LIBOR was a panel of banks' estimate of the rate at which they could borrow
unsecured for a term (one to twelve months), set each morning for the period
ahead. SOFR is a transaction-based overnight rate for repo secured by
Treasuries, compounded over the period and known at its end. LIBOR contained
bank credit risk and a term premium; SOFR contains neither. A loan that moved
to SOFR adds a spread (26.161 basis points for three-month dollar LIBOR in
the fallback) or a higher margin to keep the lender whole.
\iqlookfor{secured versus unsecured, term versus overnight, forward-looking
versus in arrears, and why the spread exists.}
\end{solution}

\begin{solution}{iq:m2:central-banks-and-the-short-rate:2}
The deposit rate binds only institutions that can use the deposit facility.
Cash-rich institutions that cannot, money-market funds, insurers, some
government agencies, lend to banks, which accept the cash only below the
deposit rate and deposit it at the central bank for the difference. An
unsecured benchmark that counts such borrowing, like the \texteuro STR, prints
below the floor; the Fed's reverse repo facility, open to money funds, puts a
second floor under that leak.
\iqlookfor{who has access to the facility, and the arbitrage that follows.}
\end{solution}

\begin{solution}{iq:m2:central-banks-and-the-short-rate:3}
Central bank: bonds up, \omterm{def:m2:central-banks-and-the-short-rate:reserves}{reserves} up. The seller's bank: \omterm{def:m2:central-banks-and-the-short-rate:reserves}{reserves} up, the
seller's deposit up. The seller: bonds down, deposit up. \omterm{def:m2:central-banks-and-the-short-rate:reserves}{Reserves} are not
``lent out'': a bank that lends creates a deposit, and the \omterm{def:m2:central-banks-and-the-short-rate:reserves}{reserves} only move
between banks; their total is set by the central bank. What QE changes is the
duration and composition the private sector holds, hence longer-term yields,
and the amount of \omterm{def:m2:central-banks-and-the-short-rate:reserves}{reserves}, hence where the overnight rate sits on the demand
curve.
\iqlookfor{the three T-accounts, and rejection of the money-multiplier
story.}
\end{solution}

\begin{solution}{iq:m2:central-banks-and-the-short-rate:4}
\omterm{def:m2:central-banks-and-the-short-rate:arrears}{Compounding in arrears} knows the rate only after the last day of the period,
too late to send an invoice and a payment on that day. A lookback uses
fixings from $p$ business days earlier with the interest period's own day
weights; an observation shift moves the whole observation period back $p$
days and uses \emph{its} weights. The shift is what a published compounded
index reproduces (a ratio of two index values), so it is simpler to verify;
the lookback matches the calendar of the loan. They give different numbers
when weekends and holidays fall differently in the two periods.
\iqlookfor{the operational reason, and a precise statement of the difference
in weights.}
\end{solution}

\begin{solution}{iq:m2:central-banks-and-the-short-rate:5}
Look at where the market overnight rate sits relative to the facility rates
and whether it responds to changes in \omterm{def:m2:central-banks-and-the-short-rate:reserves}{reserves}. In a floor it sits at or just
below the deposit rate and barely moves when \omterm{def:m2:central-banks-and-the-short-rate:reserves}{reserves} change by large
amounts (tax days, the government's account); in a corridor it sits inside
the corridor and jumps with reserve shocks and at quarter- and month-ends.
The central bank's balance sheet (\omterm{def:m2:central-banks-and-the-short-rate:reserves}{reserves} much larger than required) and its
daily operations (few or none) confirm it.
\iqlookfor{spread to the deposit rate and sensitivity to reserve supply as
the two tests.}
\end{solution}

\begin{solution}{iq:m2:central-banks-and-the-short-rate:6}
Weekends and holidays (weights of 3 or 4 days; the right holiday calendar,
which for SOFR is the US government-securities calendar); a lookback that
reaches back into the previous period; an observation shift at month-ends
where the two periods have different numbers of business days; lockouts;
missing or republished fixings; negative rates; the period's first day being a
holiday; the day-count basis by currency; rounding of the compounded rate
and of the interest amount, and of a published index if used; a period of one
day; payment dates after the period end; consistency with the published
compounded index to its rounding.
\iqlookfor{a systematic list organised by calendar, conventions, data and
rounding, and tests that compare against an independent source.}
\end{solution}
